\chapter{Introduction}
\label{ch:intro}

While testing TinyLlama on simple questions during this project, we found its answers fluent but often beside the point. Asked ``what is 2 + 3?'' in one of our demo sessions, the base model replied with a short Java program that prints 5. The answer was not wrong, but no person asking that question wants it. This gap between producing likely text and producing the answer a user actually wanted is what preference alignment tries to close. Reinforcement Learning from Human Feedback (RLHF) does this by collecting human choices between two candidate answers and using those choices to train the model. In this chapter we explain the background to RLHF, why we chose to mix two preference datasets with very different labelling styles, what problem we set out to solve, our objectives, and how the report is laid out.

\section{Background}
\label{sec:background}

\subsection{From language modelling to assistants}
GPT-3~\cite{brown2020gpt3}, Llama~2~\cite{touvron2023llama2} and TinyLlama~\cite{zhang2024tinyllama} all use the Transformer~\cite{vaswani2017attention} and are pre-trained on a very large text corpus, between a few hundred billion and a few trillion tokens. The training target is always the next token. That target teaches grammar and a great deal of factual knowledge, but nothing in it rewards following the user's request. Given a question as input, a raw pre-trained model treats it as the start of a document. It may answer, but it is just as likely to add more questions after it or carry on in a different direction, since text of all these kinds appears after questions on the web.

The usual fix is to train the model further in stages. InstructGPT~\cite{ouyang2022instructgpt} made the following three-step recipe the standard:
\begin{enumerate}
  \item \textbf{Supervised fine-tuning (SFT):} the model is trained on example conversations written to show good behaviour.
  \item \textbf{Reward modelling:} people are shown two outputs for the same prompt and pick the better one; a second network learns to predict these picks as a score.
  \item \textbf{Policy optimisation:} the language model, now called the \emph{policy}, is updated so that its outputs get higher scores, with a penalty that stops it moving too far from the model it started as.
\end{enumerate}
A striking result from that paper is that human raters preferred the 1.3B-parameter InstructGPT model to the original 175B-parameter GPT-3. A model over a hundred times smaller won because it had been aligned.

\subsection{Why pairwise preferences?}
Absolute ratings are hard to collect reliably. A response that one annotator would score 7 out of 10 might get a 5 from someone else, and the same person may not even be consistent from day to day. Picking the better of two responses is a much easier judgement, and annotators agree on it far more often. For this reason almost all RLHF datasets are built from comparisons~\cite{christiano2017deep,stiennon2020summarize,bai2022hh}. The standard statistical model for such data is the Bradley--Terry model~\cite{bradley1952rank}: every item gets a hidden score, and the chance that item A beats item B is the logistic (sigmoid) function of score(A) minus score(B). The reward model we train in Chapter~\ref{ch:method} is simply a network that outputs this hidden score for a (prompt, response) pair.

\subsection{From PPO to Direct Preference Optimization}
In the InstructGPT recipe the third step uses Proximal Policy Optimization (PPO)~\cite{schulman2017ppo}. PPO works well at scale, but it is expensive. Four models have to be held on the GPU together: the policy being trained, a frozen copy of it used as the reference, the reward model, and a value network that estimates expected reward. New responses must be generated at every training step, and the method is known to be hard to tune. On the free Kaggle T4 GPUs available to us, which have about 16\,GB of memory each and no BF16 support, this was not realistic.

Direct Preference Optimization (DPO)~\cite{rafailov2023dpo} removes most of this cost. Its authors showed that the best policy under the PPO objective can be written down in closed form, and that training towards it reduces to a logistic loss on the stored preference pairs. No separate reward model or value network is needed during training and nothing is generated on the fly, so DPO trains much like ordinary fine-tuning with only the policy and the reference in memory. We combined DPO with 4-bit quantisation and LoRA adapters~\cite{hu2022lora,dettmers2023qlora}. With these, the 1.1B-parameter TinyLlama took only about 1.1\,GB of GPU memory once loaded, leaving plenty of room for training.

\section{Motivation}
\label{sec:motivation}

\subsection{The single-dataset blind spot}
Most RLHF tutorials and student projects we looked at train on one dataset, usually Anthropic's HH-RLHF~\cite{bai2022hh}, and then test the reward model on pairs held back from that same dataset. A good score in that setting says less than it seems to. Every preference dataset carries the habits of the people and process that produced it: who the annotators were, what instructions they followed, what the interface looked like and how the conversations were formatted. A reward model can do well on test pairs from the same source simply by picking up these habits, without learning much about what makes an answer good in general.

HH-RLHF is a clear example. All of its 43,835 helpful-base training pairs are dialogues in a fixed \texttt{Human:}/\texttt{Assistant:} format, labelled by paid contractors working from a helpfulness guide. A model trained only on this data could end up rewarding the style of HH-RLHF's preferred answers rather than their quality. LMSYS Chatbot Arena~\cite{chiang2024arena} is different in almost every respect. Its votes come from anonymous members of the public comparing answers from real chatbots, its answers are longer, and in the copy we used 17,761 of 57,477 comparisons (about 31\%) ended in a tie. A model trained only on Arena votes could instead learn whatever length or formatting that crowd tended to favour.

\subsection{Why combine expert and crowd labels}
Putting both datasets into one training set does two jobs for us. The obvious one is variety: the reward model sees the careful, guideline-based judgements of HH-RLHF and also the much wider range of real questions that people type into Chatbot Arena. The less obvious job turned out to be more useful. Because every example keeps a tag saying which dataset it came from, we can test the reward model on HH-RLHF pairs and on LMSYS pairs separately. If it does well on one and badly on the other, it has probably learned that dataset's habits. If the two scores are close, it has learned something that carries over between the two styles of labelling. We apply the same split to the DPO model's results.

\subsection{Working on free hardware}
We had no access to a GPU cluster, so a second aim was to see how much of an RLHF study can be done with free tools. Everything in this report was trained on Kaggle's free notebooks (two NVIDIA T4 cards with about 16\,GB each), and a smaller version was later trained on an ordinary 8\,GB Apple M1 laptop. The first full training run takes about two and a half to three hours. After that the trained adapters are saved, and every later evaluation or demonstration starts in under a minute.

\section{Problem Statement}
\label{sec:problem}

\begin{quote}
\itshape
Build and test a preference-alignment pipeline that (i) merges expert-labelled Anthropic HH-RLHF data and crowd-voted LMSYS Chatbot Arena data into one balanced training set in which every example is tagged with its source; (ii) trains a reward model on this set and measures its accuracy on each source separately, to detect whether it has latched onto one source; (iii) aligns a small open chat model with Direct Preference Optimization within the memory of a free 16\,GB GPU; and (iv) checks, again per source, whether the aligned model is preferred to the original on prompts it has not seen.
\end{quote}

Three practical problems sit inside this statement. The \emph{data} problem is that the two datasets store conversations in completely different ways (one long dialogue string versus JSON lists with three winner columns), have different answer lengths, and in LMSYS almost a third of the rows are ties. The \emph{compute} problem is fitting a 1.1B-parameter policy, its reference copy and a 184M-parameter reward model into 16\,GB. The \emph{evaluation} problem is measuring improvement without paying for human raters, while being clear about what an automatic judge can and cannot tell us.

\section{Objectives}
\label{sec:objectives}

\begin{enumerate}[label=\textbf{O\arabic*.}]
  \item Write loaders that bring HH-RLHF (helpful-base) and LMSYS Chatbot Arena into one \{prompt, chosen, rejected, source\} format, remove ties and broken rows, apply the same length rule to both, and produce train, validation and test splits that each contain both sources.
  \item Train a reward model on DeBERTa-v3-base~\cite{he2023debertav3} with LoRA adapters and the Bradley--Terry loss, and save it so it does not need retraining.
  \item Measure the reward model's pairwise accuracy and average score gap on held-out pairs, overall and per source, and treat a difference of under 10 percentage points between the sources as ``balanced''.
  \item Align TinyLlama-1.1B-Chat~\cite{zhang2024tinyllama} with DPO using 4-bit weights, LoRA, gradient checkpointing and an 8-bit paged optimiser so that it trains on one T4.
  \item Compare the original and aligned models on five fixed demo prompts and on 50 unseen prompts (25 from each source), reporting win rate and average score change per source.
  \item Package the work as a reusable system: a command-line pipeline that runs each stage, and a local web application that shows the two models' answers side by side with their scores.
  \item Discuss the results critically, including the weaknesses of using the reward model as the judge, and list concrete next steps.
\end{enumerate}

\section{Scope of the Project}
\label{sec:scope}

\paragraph{Included.} Stages two and three of the RLHF recipe (reward modelling and preference optimisation) for English helpfulness: data preparation for two public datasets, parameter-efficient training of both models, automatic per-source evaluation, two full training configurations on Kaggle and one reduced configuration on a laptop, and a web application for live comparison.

\paragraph{Not included.} Pre-training, a separate supervised fine-tuning stage (we start from a model that already has chat training), PPO, and explicit measurement of harmlessness or factual correctness. Large human studies and evaluation by a strong external model are left for future work. Both of our models are small (184M and 1.1B parameters), so we make no claims about behaviour at 7B parameters or more.

\paragraph{Assumptions and limits.}
\begin{itemize}
  \item Only free hardware is used: Kaggle T4 GPUs, which cannot do BF16 arithmetic and stop after a session time limit, and an 8\,GB Apple M1 laptop.
  \item The ``chosen'' and ``rejected'' labels in both datasets are taken as correct, even though both contain noise.
  \item Text fields are cut to 1,500 characters and model inputs to 512 tokens (fewer on the laptop) to save memory, which can remove useful context from long conversations.
  \item Kaggle deletes \emph{/kaggle/working} between sessions, so trained adapters must be saved as notebook output and attached again as input to be reused.
\end{itemize}

\section{Brief Methodology}
\label{sec:brief-method}

Figure~\ref{fig:intro-overview} gives the outline. The two datasets are converted to one format, cleaned and balanced. A reward model is trained on the training split. The same split supplies preference pairs for DPO. Both models are then tested per source on held-out data, and the trained models are served through a web application. Chapter~\ref{ch:method} describes the design and Chapter~\ref{ch:impl} the implementation.

\begin{figure}[H]
\centering
\begin{tikzpicture}[
  blk/.style={draw, rounded corners=3pt, align=center, minimum height=1.35cm, text width=2.5cm, font=\footnotesize, drop shadow={opacity=0.15}},
  arr/.style={-{Stealth[length=2.4mm]}, thick, color=black!70}]
\node[blk, fill=boxA] (hh) at (0,1.0) {Anthropic HH-RLHF\\(expert labels)};
\node[blk, fill=boxA] (lm) at (0,-1.0) {LMSYS Arena\\(crowd votes)};
\node[blk, fill=boxB] (comb) at (3.5,0) {Normalise, filter, balance, stratify (source-tagged)};
\node[blk, fill=boxC] (rm) at (7.0,1.0) {Reward model\\DeBERTa-v3 + LoRA};
\node[blk, fill=boxD] (dpo) at (7.0,-1.0) {DPO alignment\\TinyLlama-1.1B};
\node[blk, fill=boxE] (ev) at (10.5,0) {Per-source evaluation (HH / LMSYS)};
\draw[arr] (hh.east) -- (comb.north west);
\draw[arr] (lm.east) -- (comb.south west);
\draw[arr] (comb.north east) -- (rm.west);
\draw[arr] (comb.south east) -- (dpo.west);
\draw[arr] (rm.east) -- (ev.north west);
\draw[arr] (dpo.east) -- (ev.south west);
\draw[arr, dashed] (rm.south) -- node[right, font=\scriptsize]{judge} (dpo.north);
\end{tikzpicture}
\caption[Pipeline overview]{High-level overview of the proposed cross-distribution RLHF pipeline.}
\label{fig:intro-overview}
\end{figure}

\section{Contributions}
\label{sec:contrib}

\begin{itemize}
  \item \textbf{One tagged dataset from two very different sources.} A parser that turns HH-RLHF dialogues and Arena JSON records into one format, removes ties and ambiguous rows, shortens rather than discards long examples, and splits each source separately so both appear in every split.
  \item \textbf{Per-source evaluation throughout.} Reward-model accuracy, score gap and DPO win rate are all reported separately for HH-RLHF and LMSYS. This made visible an effect that the overall numbers hid.
  \item \textbf{A complete pipeline on free hardware.} DeBERTa-v3-base and TinyLlama-1.1B-Chat trained end to end in about 2.5--3 hours on Kaggle T4 GPUs, with saving and reloading so that later runs start in seconds.
  \item \textbf{A working system.} The notebook code reorganised into a Python package with one command per stage, a reduced configuration that trains on an 8\,GB laptop, and a web application that compares the two models live and shows the evaluation results.
  \item \textbf{Measured results reported honestly.} Two full Kaggle configurations and one laptop configuration are compared. Measured numbers are kept apart from the higher figures in an earlier paper draft, and we explain the difference, including a case where the reward model rewarded a wrong answer.
  \item \textbf{Practical fixes for T4 and M1 hardware.} FP32 for DeBERTa-v3, FP16 compute for the 4-bit policy, saving the reward model's pooler, correct device placement on two GPUs, leakage-free splitting and lower-memory DPO on the laptop, all documented in Chapter~\ref{ch:impl}.
\end{itemize}

\section{Organisation of the Report}
\label{sec:organisation}

\begin{description}[style=nextline, leftmargin=2.2cm, font=\normalfont\bfseries]
  \item[Chapter 2] reviews the papers on RLHF, reward models, DPO and its variants, parameter-efficient training, preference datasets and evaluation, and lists the gaps we address.
  \item[Chapter 3] presents the design: the architecture, data-flow diagrams, the mathematics of the reward model, DPO and LoRA, the two models, the datasets, the metrics and the tools.
  \item[Chapter 4] describes the implementation module by module with code and algorithms, the packaged system and web application, and the problems we solved along the way.
  \item[Chapter 5] gives the experimental set-up and results for the two Kaggle configurations and the laptop run, qualitative examples, the web application in use, and a discussion that also covers the figures from our earlier paper draft, clearly marked.
  \item[Chapter 6] summarises what was achieved, the limitations and possible next steps.
\end{description}
The references and appendices, with full code listings, sample outputs and instructions for re-running the work, follow the chapters.
